% Promoted exact kernel calculation after full root and companion audits. Candidate retained; no restricted-cup lower bound asserted.
\subsection{The Frobenius-trivial extension and the residual-boundary kernel}

\begin{lemma}[Frobenius of the extension class and its symmetric square]
\label{lem:DE-A3-unipotent-Frobenius}
\status{Proved; complete independent root and companion audits passed.}
In the setting of the explicit unipotent quotient lemma, the extension
class $\gamma$ satisfies $F\gamma=0$ in $H^1(T,\mathcal O_T)$.
If its actual local polarization index is zero, then
\[
 \operatorname{Sym}^2G\simeq U\oplus\mathcal O_T,
 \qquad(h^0,h^1,h^2)(\operatorname{Sym}^2G)=(2,4,2),
 \quad G=\mathfrak g^\vee\simeq UP.
 \tag{DE-SK1}
\]
If the actual local index is positive, the corresponding bundle is
$\operatorname{Sym}^2G\simeq\mathcal O_T^3$, with cohomology dimensions
$(3,6,3)$. The displayed splittings need not be canonical.
\end{lemma}
\begin{proof}
Keep the actual local functions from (DE-UB1). Since
$h_j-h_i=c_{ij}^2$, the differentials $dh_i$ glue to a global one-form
$\theta$ on $T$. It is nonzero: the finite cover is locally
$z_i^2=h_i$, and smoothness of its total space forces $dh_i$ to be
nonzero. The extension formula gives, without a choice of square root,
\[
 d\gamma_{ij}=\gamma_{ij}^2\theta.
\]
The map $d:H^1(\mathcal O_T)\to H^1(\Omega_T^1)$ is zero by
Hodge--de Rham degeneration for abelian varieties
\cite[Corollary~7.29 and its full proof, pp.~108--109]
{EdixhovenVanDerGeerMoonenAV}.
Moreover $\Omega_T^1=H^0(\Omega_T^1)\otimes\mathcal O_T$, so cup
product with the nonzero invariant form $\theta$ injects
$H^1(\mathcal O_T)$ into $H^1(\Omega_T^1)$.
The displayed cocycle identity therefore gives $F\gamma=0$.
This argument also covers the split case.

Choose an affine cover and functions $f_i$ with
$\gamma_{ij}^2=f_j-f_i$. In a frame of $U$ with transition
$\left(\begin{smallmatrix}1&\gamma_{ij}\\0&1\end{smallmatrix}\right)$,
the symmetric-square frame $e_1^2,e_1e_2,e_2^2$ has transition
\[
 \begin{pmatrix}1&\gamma_{ij}&\gamma_{ij}^2\\0&1&0\\0&0&1\end{pmatrix}.
\]
Replacing $e_2^2$ by $e_2^2+f_i e_1^2$ removes its last off-diagonal
entry. This proves $\operatorname{Sym}^2U\simeq U\oplus\mathcal O_T$.
Because $P^2\simeq\mathcal O_T$ and $U$ is self-dual, the same is true
of $\operatorname{Sym}^2G$.
For nonzero $\gamma$, the extension $0\to\mathcal O_T\to U\to
\mathcal O_T\to0$ has $h^0(U)=h^2(U)=1$ and $h^1(U)=2$:
its connecting maps are cup product with $\gamma$, of rank one in
degrees zero and one, using the perfect alternating cup pairing on
$H^1(\mathcal O_T)$ \cite[Proposition~7.27, p.~108]
{EdixhovenVanDerGeerMoonenAV}. This gives (DE-SK1).
When $\gamma=0$, $U=\mathcal O_T^2$, which gives the final case.
\end{proof}

\begin{lemma}[The additional twisted-pencil contribution to the boundary kernel]
\label{lem:DE-A3-residual-boundary-kernel}
\status{Proved; complete independent root and companion audits passed.}
At a sharp source retain $\mathcal L=LP$, $E=Q^\vee$, $G=\mathfrak g^\vee$,
and $R=\omega_D A^{-1}$ from (DE-RB1)--(DE-RB2).
Let $\mathcal J_s$ be the image of $E\otimes G\to\operatorname{Sym}^2G$.
There is an exact sequence
\[
 0\longrightarrow H^0(T,\operatorname{Sym}^2G)
 \longrightarrow\ker\partial_s
 \longrightarrow
 \ker\!\left[H^1(T,E\otimes G)
       \longrightarrow H^1(T,\operatorname{Sym}^2G)\right]
 \longrightarrow0.
 \tag{DE-SK2}
\]
The first arrow is restriction along the quotient in (DE-RB2).
Tensoring the squared-pencil sequence by $G$ identifies the remaining
domain exactly as
\[
 H^1(T,E\otimes G)\simeq H^0(D,A\otimes G|_D).
 \tag{DE-SK3}
\]
In particular one may not discard this domain using acyclicity of $G$.
Indeed the known nonzero norm section $\sigma\in H^0(D,RP)$ gives
$h^0(D,AP)=h^0(D,RP)\ge1$ by curve duality and $\chi(AP)=0$;
the subline $P\subset G$ then gives
$h^0(D,A\otimes G|_D)\ge1$.
\end{lemma}
\begin{proof}
The squared-pencil sequence is $0\to E\to\mathcal O_T^2\to i_*A\to0$.
Tensor it by the locally free $G$. Because $G$ is acyclic, its long
cohomology sequence proves (DE-SK3), and its initial injection gives
$H^0(E\otimes G)=0$.
On the abelian surface, $H^1(\mathcal L^{-1})=0$ by Serre duality
and vanishing for ample line bundles. The first short exact sequence
in (DE-RB2),
\[
0\to\mathcal L^{-1}\to E\otimes G\to\mathcal J_s\to0,
\]
therefore gives $H^0(\mathcal J_s)=0$ and an injection
$H^1(E\otimes G)\hookrightarrow H^1(\mathcal J_s)$.
The second short exact sequence is
$0\to\mathcal J_s\to\operatorname{Sym}^2G\to i_*R^2\to0$.
Its boundary has kernel $H^0(\operatorname{Sym}^2G)$ and image the
kernel of $H^1(\mathcal J_s)\to H^1(\operatorname{Sym}^2G)$.
Composing with the first boundary gives $\partial_s$. Taking the
preimage of $H^1(E\otimes G)$ proves (DE-SK2), including surjectivity
onto the displayed kernel. All these sequences apply equally to a
nonreduced Cartier divisor.

Finally $P^2=\mathcal O_T$ and $\omega_D A^{-1}=R$ give
$H^1(AP)^*=H^0(RP)$. Since $\deg A=g(D)-1$, one has $\chi(AP)=0$.
The regular norm section already constructed in (DE-TP2) is nonzero,
so both spaces have positive dimension. Restricting the actual
subbundle $P\subset G$ and tensoring by $A$ proves the last assertion.
\end{proof}

\begin{remark}[What remains after the cohomology calculation]
\label{rem:DE-A3-symmetric-boundary-kernel-scope}
\status{Proved; complete independent root and companion audits passed.}
The computed dimensions of $H^0(\operatorname{Sym}^2G)$ control only
one explicit subspace of $\ker\partial_s$. Formula (DE-SK2) retains
the additional twisted-pencil map, and neither formula controls the
intersection with the three-dimensional residual products $W_s$ or
the subsequent restriction to the actual direction space $\Delta_s$.
Thus no lower bound on $\dim\operatorname{im}\Psi_s$ follows here.
The local symmetric-square change of frame is among the seventeen
universal identities in the retained unipotent transition certificate;
no additional numerical experiment or finite-sample inference is used.
\end{remark}
